%=======================================================================
% setup_teaching.tex -- the entire preamble for this deck.
%
% Session 3, "AI Tools: Introduction to Claude Code"
% Yale MAM Orientation BootCamp
%
% slide.tex inputs this file and nothing else. Everything the deck needs
% is here: theme, colours, text macros, code blocks, frame furniture.
%
% Build with pdflatex. fontenc[T1] + inputenc[utf8] is a pdflatex setup;
% xelatex and lualatex want fontspec instead.
%
% THE ONE TRAP
% ------------
% Any frame containing a prompt / pycode / shellout block MUST be
% declared
%
%     \begin{frame}[fragile]{Title}
%
% Without [fragile], beamer scans the frame body before listings can
% claim it, the \end{...} is never seen, and the run dies pages later
% with "Text dropped after begin of listing" followed by an emergency
% stop pointing at the end of the file rather than at your frame. This
% is the most likely way to break this deck.
%
% House style: ASCII punctuation. Write "--" for a dash and "'" for an
% apostrophe. Both render identically to their Unicode forms and survive
% being copied into a terminal. UTF-8 compiles fine, so an accented
% author name is not a problem.
%=======================================================================


%=======================================================================
% 1. Theme and packages
%=======================================================================
% Add a package here only when a frame needs it, with a comment saying
% which one. The list is short on purpose.

\usetheme{Szeged}
\usecolortheme{seahorse}

\usepackage[T1]{fontenc}
\usepackage[utf8]{inputenc}
\usepackage[english]{babel}

\usepackage{lmodern}          % scalable Latin Modern. Supplies the
                              % monospace face the listings blocks use:
                              % lato has no mono, and without lmodern
                              % the code blocks fall back to bitmap cmtt
\usepackage[default]{lato}    % text font, matching the research decks

\usepackage{xcolor}
\usepackage{graphicx}
\usepackage{booktabs}         % \toprule \midrule \bottomrule
\usepackage{colortbl}         % \rowcolor, to tint one row of a table
\usepackage{amsmath}
\usepackage{listings}
\usepackage{tikz}             % box-and-arrow diagrams; see \flowbox below
\usetikzlibrary{arrows.meta,positioning}
\usepackage{hyperref}
\usepackage{pgfpages}         % required by \setbeameroption below

% Speaker notes. This deck ships with notes hidden. To write them, put
% \note{...} after a frame and switch the option. Presenting from a
% laptop with an external screen wants the third form.
\setbeameroption{hide notes}
%\setbeameroption{show only notes}
%\setbeameroption{show notes on second screen=right}


%=======================================================================
% 2. Colours and text macros
%=======================================================================

\definecolor{yaleblue}{HTML}{00356b}
\definecolor{ylb}{HTML}{286dc0}
\definecolor{yo}{HTML}{bd5319}
\definecolor{lightgray}{HTML}{d3d3d3}

% Emphasis, in descending order of how much attention it claims:
%   \ybb  Yale blue bold  -- main emphasis. The finding on a slide.
%   \ylbb light blue bold -- secondary emphasis.
%   \yob  orange bold     -- warnings, and the thing that broke.
\newcommand\yb[1]{{\color{yaleblue} #1}}
\newcommand\ybb[1]{{\color{yaleblue} \textbf{#1}}}
\newcommand\ylb[1]{{\color{ylb} #1}}
\newcommand\ylbb[1]{{\color{ylb} \textbf{#1}}}
\newcommand\yo[1]{{\color{yo} #1}}
\newcommand\yob[1]{{\color{yo} \textbf{#1}}}

% \blue, \red and \green are left undefined on purpose. The MAM course
% decks in slides/misc define those names with different colours, so a
% frame copied in from one of them fails loudly with "Undefined control
% sequence" rather than quietly drawing the wrong colour. Use the Yale
% macros above.


%=======================================================================
% 3. Beamer appearance
%=======================================================================

\setbeamercolor{structure}{fg=white}
\setbeamercolor{title}{fg=yaleblue,bg=white}
\setbeamercolor{frametitle}{fg=yaleblue,bg=white}
\setbeamercolor{button}{bg=yaleblue!10,fg=yaleblue}

\setbeamerfont{frametitle}{size=\fontsize{14pt}{12pt}}
\setbeamertemplate{frametitle}[default][center]
\addtobeamertemplate{frametitle}{\vspace*{0cm}}{\vspace*{-0.5cm}}

\setbeamertemplate{itemize items}[default]
\setbeamercolor{itemize item}{fg=ylb,bg=white}
\setbeamercolor{itemize subitem}{fg=ylb,bg=white}
\setbeamercolor{itemize subsubitem}{fg=ylb,bg=white}
\setbeamertemplate{itemize subitem}[triangle]
\setbeamertemplate{itemize subsubitem}[triangle]

\setbeamertemplate{enumerate items}[default]
\setbeamercolor{enumerate item}{fg=ylb,bg=white}
\setbeamercolor{enumerate subitem}{fg=ylb,bg=white}

\setbeamertemplate{navigation symbols}{}
\setbeamersize{text margin left=0.5cm,text margin right=0.5cm}

% Tighter display maths.
\makeatletter
\g@addto@macro\normalsize{
  \setlength\abovedisplayskip{5pt}
  \setlength\belowdisplayskip{5pt}
  \setlength\abovedisplayshortskip{5pt}
  \setlength\belowdisplayshortskip{5pt}
}
\makeatother

\setbeamertemplate{section page}{
  \begin{centering}
  \fontsize{20}{16} {\color{yaleblue} \textsc{\insertsection}} \par
  \end{centering}
}

%----------------------------------------------------------------------
% 3a. No navigation bar across the top
%----------------------------------------------------------------------
% Szeged puts a miniature table of contents in the headline: a row of
% section names with a progress dot per frame. It is blanked here. Over
% 86 frames the dots compress to unreadable pinheads, the bar costs
% about a centimetre of vertical space on every slide, and none of it
% reads from the back of a lecture hall.
%
% To put it back, delete this one line.
\setbeamertemplate{headline}{}

% The footline carries the position information instead: short title on
% the left, institute and then the frame number on the right.
% Keep the frame number if you restyle this -- it is what lets someone
% in row 12 ask about "slide 47".
%
% The left half prints \insertshorttitle (the bracketed argument of
% \title) and the right half \insertshortinstitute (the bracketed
% argument of \institute), so both are set in slide.tex, not here.
\setbeamercolor{footline}{fg=yaleblue}
\setbeamerfont{footline}{size=\fontsize{6.5pt}{8pt}}
\setbeamertemplate{footline}{%
  \hbox{%
    \begin{beamercolorbox}[wd=.55\paperwidth,ht=2.4ex,dp=1.2ex,leftskip=.3cm]%
      {footline}\usebeamerfont{footline}\insertshorttitle
    \end{beamercolorbox}%
    \begin{beamercolorbox}[wd=.45\paperwidth,ht=2.4ex,dp=1.2ex,rightskip=.3cm]%
      {footline}\usebeamerfont{footline}\hfill
      \insertshortinstitute{} \quad \insertframenumber{} / \inserttotalframenumber
    \end{beamercolorbox}%
  }%
  \vskip0pt%
}


%=======================================================================
% 4. Code, prompts and terminal output
%=======================================================================
% White monospace on a cobalt block. Remember [fragile] on any frame
% that uses these; see the file header.

\definecolor{cobalt}{RGB}{0, 34, 64}
\definecolor{bluecomments}{RGB}{0, 103, 231}
\definecolor{colnum}{RGB}{251, 94, 133}
\definecolor{stringgreen}{RGB}{35, 185, 23}

% LST1: the default. "language=" is empty on purpose. This deck shows
% Python, shell and plain English prompts in the same style of block,
% and keyword highlighting applied to an English sentence looks like a
% rendering bug.
\lstdefinestyle{LST1}{
    language=,
    basicstyle=\scriptsize\ttfamily\color{white},
    commentstyle=\color{bluecomments},
    stringstyle=\color{stringgreen},
    numberstyle=\color{colnum},
    literate={~}{$\sim$}{1},
    backgroundcolor=\color{cobalt},
    showspaces=false,
    showstringspaces=false,
    linewidth=\linewidth,
    breaklines=true,
    breakatwhitespace=true,
    columns=fullflexible,
    keepspaces=true,
    escapeinside={(*@}{@*)}
}

% LST_fn: the same block one size up, for a short prompt that is the
% entire point of its slide and must read from the back row.
\lstdefinestyle{LST_fn}{
    language=,
    basicstyle=\footnotesize\ttfamily\color{white},
    commentstyle=\color{bluecomments},
    stringstyle=\color{stringgreen},
    numberstyle=\color{colnum},
    literate={~}{$\sim$}{1},
    backgroundcolor=\color{cobalt},
    showspaces=false,
    showstringspaces=false,
    linewidth=\linewidth,
    breaklines=true,
    breakatwhitespace=true,
    columns=fullflexible,
    keepspaces=true,
    escapeinside={(*@}{@*)}
}

% LSTpy: Python, for the few slides showing real source.
\lstdefinestyle{LSTpy}{
    language=Python,
    basicstyle=\scriptsize\ttfamily\color{white},
    keywordstyle=\color{ylb}\bfseries,
    commentstyle=\color{bluecomments},
    stringstyle=\color{stringgreen},
    backgroundcolor=\color{cobalt},
    showspaces=false,
    showstringspaces=false,
    linewidth=\linewidth,
    breaklines=true,
    breakatwhitespace=true,
    columns=fullflexible,
    keepspaces=true,
    escapeinside={(*@}{@*)}
}

\lstset{style=LST1}

% \begin{prompt} ... \end{prompt}
% A prompt as it was actually run. Every prompt in this deck is quoted
% verbatim; the same text is in ../project/prompts/. Do not paraphrase
% one on a slide: the point of showing it is that those exact words did
% the work.
\lstnewenvironment{prompt}[1][]
  {\lstset{style=LST_fn,#1}}
  {}

% \begin{pycode}   -- Python source
% \begin{shellout} -- terminal or agent output
%
% Named pycode and shellout because beamer already claims the names
% "code" and "output"; \lstnewenvironment refuses to clobber an
% existing environment.
\lstnewenvironment{pycode}[1][]
  {\lstset{style=LSTpy,#1}}
  {}
\lstnewenvironment{shellout}[1][]
  {\lstset{style=LST1,#1}}
  {}


%=======================================================================
% 5. Frame furniture
%=======================================================================

% Looser lists. Prefer wideitemize on lecture frames; beamer's default
% spacing packs bullets too tightly to scan from a distance.
\newenvironment{wideitemize}{\itemize\addtolength{\itemsep}{8pt}}{\enditemize}
\newenvironment{wideenumerate}{\enumerate\addtolength{\itemsep}{8pt}}{\endenumerate}

% A part divider: nothing but the part title, large, Yale blue, centred
% on a white slide.  Used once at the top of each of the four parts so
% the room always knows where it is in three hours.
%
% The body is centred both ways, so write only the title inside:
%
%     \begin{transitionframe}
%     Part 2: Building a private credit dataset
%     \end{transitionframe}
%
% Sizing is set here (\Large, bold) -- do not re-size inside the body,
% or the four dividers stop matching each other.
\newenvironment{transitionframe}{
  \setbeamercolor{background canvas}{bg=white}
  \setbeamercolor{normal text}{fg=yaleblue}
  \usebeamercolor[fg]{normal text}
  \centering
  \begin{frame}[plain]
  \vfill
  \leftskip=1.2cm \rightskip=1.2cm \parindent=0pt
  \centering
  \huge\bfseries}{
  \par\vfill
  \end{frame}
}

% \takeaway{...} -- the one sentence a slide exists to deliver, as a
% tinted full-width band. At most one per frame, and deliberately not on
% every frame: if every slide has one, none of them land.
\newcommand{\takeaway}[1]{%
  \vspace{0.3em}%
  \begin{beamercolorbox}[wd=\textwidth,sep=6pt,rounded=true]{button}%
  \small #1%
  \end{beamercolorbox}%
}

% \source{...} -- provenance line at the foot of the frame, deliberately
% quiet. Every frame stating a measured number carries one, naming the
% memo the number came from. This is not decoration: the deck argues
% that unsourced numbers are the problem it is teaching about.
\newcommand{\source}[1]{%
  \vfill{\scriptsize\color{gray} #1\par}}

% \capture{height fraction}{what to shoot} -- a dashed grey box standing in
% for a screenshot that has not been taken yet.  It renders in the PDF on
% purpose, so a missing capture is visible when the deck is reviewed rather
% than silently absent.  To fill it, replace the whole \capture call with
%     \begin{center}\includegraphics[width=...]{file}\end{center}
% Nothing else on the frame needs to change.
\newcommand{\capture}[2]{%
  \begin{center}
  \begin{tikzpicture}
    \node[draw=lightgray, dashed, line width=0.8pt, rounded corners=3pt,
          fill=lightgray!12, text=gray, align=center, font=\small,
          text width=0.78\textwidth, minimum height=#1\textheight]
         {capture needed: #2};
  \end{tikzpicture}
  \end{center}}

%----------------------------------------------------------------------
% 5a. Box-and-arrow diagrams
%----------------------------------------------------------------------
% Three tikz node styles for flow diagrams:
%   flowbox       a neutral step, or one the tool takes
%   flowbox_us   a step you take -- filled Yale blue
%   flowbox_warn  the thing that went wrong -- orange
% The deck uses the colour difference to carry an argument, so pick the
% style by meaning, not by position. Draw arrows with [flowarrow].
%
%   \begin{tikzpicture}[node distance=4mm]
%     \node[flowbox_us] (a) {instruct};
%     \node[flowbox, right=of a] (b) {it acts};
%     \draw[flowarrow] (a) -- (b);
%   \end{tikzpicture}
%
% Box width is fixed so a row of them reads as one strip. Change
% text width= here, not per diagram.
\tikzset{
  flowbox/.style={
    draw=yaleblue!40, line width=0.6pt, rounded corners=2pt,
    fill=yaleblue!5, text=yaleblue!80!black,
    align=center, font=\small,
    text width=2.55cm, inner sep=5pt, minimum height=1.15cm},
  flowbox_us/.style={
    flowbox, fill=yaleblue, draw=yaleblue, text=white, font=\small\bfseries},
  flowbox_warn/.style={
    flowbox, fill=yo!8, draw=yo!50, text=yo},
  flowarrow/.style={
    -{Stealth[length=2.4mm]}, line width=0.9pt, color=yaleblue!70},
  flowlabel/.style={font=\scriptsize, text=yo, align=center},
}

\newcommand{\tabfont}{\fontsize{9}{10}\selectfont}  % shrink a wide table
\newcommand{\file}[1]{{\ttfamily\small #1}}          % file and field names
\newcommand{\vs}{\vspace{0.5em}}
